
\subsection{Phase Space Representation of an Optical Beam}


In paraxial (ray) optics, a ray propagating along the $y$ direction can be studied in a transverse direction ($x$ or $z$, we use $x$) by two coordinates:
\begin{equation}
  \mathbf{r} = \begin{pmatrix} x \\ x' \end{pmatrix},
\end{equation}
where
\begin{itemize}
  \item $x$ is the transverse height of the ray at the current plane, and
  \item $x' = dx/dy$ is the paraxial angle (slope) with respect to the optical axis.
\end{itemize}

A \emph{beam} is a bundle of many rays; it is represented as a cloud of points
$(x_i,\, x'_i)$ in the phase plane.


A ray travelling a distance $d$ through a uniform medium of refractive index $n$
is transformed by
\begin{equation}
  \begin{pmatrix} x \\ x' \end{pmatrix}_{\!\text{out}}
  =
  \underbrace{\begin{pmatrix} 1 & d/n \\ 0 & 1 \end{pmatrix}}_{M_{\text{drift}}}
  \begin{pmatrix} x \\ x' \end{pmatrix}_{\!\text{in}},
\end{equation}
so explicitly
\begin{align}
  x_{\text{out}}  &= x_{\text{in}} + \frac{d}{n}\,x'_{\text{in}}, \\[4pt]
  x'_{\text{out}} &= x'_{\text{in}}.
\end{align}
In phase space this is a \textbf{horizontal shear}: the angle is unchanged, but
each ray's position shifts by an amount proportional to its slope and to the
propagation distance.  Geometrically, every point in the phase-plane cloud slides
to the right by $d\,x'/n$, so the cloud as a whole shears.

Two important limiting cases:
\begin{itemize}
  \item \textbf{Point source} at $x=0$: the cloud is a vertical line that shears
        into a tilted line --- the beam diverges spatially.
  \item \textbf{Collimated beam} ($x'=0$ for all rays): the cloud is a horizontal
        line and is unaffected by the drift --- the beam width is preserved.
\end{itemize}

% -------------------------------------------------------------------
\paragraph{The ABCD Ray-Transfer Matrix}

Any \emph{paraxial} optical element can be described by a $2\times 2$
\emph{ray-transfer} (or ABCD) matrix
\begin{equation}
  M = \begin{pmatrix} A & B \\ C & D \end{pmatrix}, \qquad
  \begin{pmatrix} x \\ x' \end{pmatrix}_{\!\text{out}}
  = M \begin{pmatrix} x \\ x' \end{pmatrix}_{\!\text{in}}.
\end{equation}
Table~\ref{tab:abcd} lists the matrices for the most common elements.

\begin{table}[h]
  \centering
  \caption{ABCD matrices for common paraxial optical elements.}
  \label{tab:abcd}
  \begin{tabular}{lcccc}
    \toprule
    Element & $A$ & $B$ & $C$ & $D$ \\
    \midrule
    Free propagation (distance $d$, index $n$) & $1$ & $d/n$ & $0$ & $1$ \\
    Thin lens (focal length $f$)               & $1$ & $0$   & $-1/f$ & $1$ \\
    Flat interface ($n_1 \to n_2$)             & $1$ & $0$   & $0$ & $n_1/n_2$ \\
    Curved mirror (radius $R$)                 & $1$ & $0$   & $-2/R$ & $1$ \\
    \bottomrule
  \end{tabular}
\end{table}

% The determinant of $M$ satisfies
% \begin{equation}
%   \det M = AD - BC = \frac{n_1}{n_2},
% \end{equation}
% and equals unity when the ray travels through a single medium.  This is the
% matrix statement of \emph{Liouville's theorem} (see Section~\ref{sec:etendue}).

% For a \emph{compound} system of elements $M_1, M_2, \ldots, M_k$ (in order of
% traversal), the overall matrix is
% \begin{equation}
%   M_{\text{sys}} = M_k \cdots M_2\, M_1.
% \end{equation}
% Note the \emph{reverse} ordering: the first element encountered by the ray is
% rightmost in the product.

% % -------------------------------------------------------------------
% \subsubsection{\'Etendue and the Optical Invariant}
% \label{sec:etendue}

% The phase-space \emph{area} occupied by the beam is proportional to the
% \textbf{\'etendue} (or optical extent):
% \begin{equation}
%   G = n \cdot \Delta x \cdot \Delta\theta,
% \end{equation}
% where $\Delta x$ and $\Delta\theta$ are the spatial and angular extents of the
% beam.  Through any lossless, stigmatic optical system the \'etendue is conserved:
% \begin{equation}
%   n_1\,x_1\,\theta_1 = n_2\,x_2\,\theta_2.
% \end{equation}
% This is the \textbf{Lagrange--Helmholtz invariant}.  Its phase-space interpretation
% is that a symplectic (area-preserving) transfer matrix can \emph{deform} the
% phase-space cloud --- stretch, rotate, shear --- but cannot change its area.
% Consequently, it is impossible to reduce the beam size and the beam divergence
% simultaneously: reducing one necessarily increases the other.

% % -------------------------------------------------------------------
% \subsubsection{Gaussian Beams and the \texorpdfstring{$q$}{q}-Parameter}

% For a \textbf{Gaussian beam} the phase-space description is most elegantly
% expressed through the \emph{complex beam parameter}
% \begin{equation}
%   q(z) = z - z_0 + i\,z_R,
% \end{equation}
% where $z_0$ is the longitudinal position of the beam waist and
% $z_R = \pi w_0^2/\lambda$ is the \emph{Rayleigh range} ($w_0$ is the $1/e^2$
% intensity half-width at the waist).  The beam width $w(z)$ and wavefront radius
% of curvature $R(z)$ are recovered from
% \begin{equation}
%   \frac{1}{q(z)} = \frac{1}{R(z)} - i\,\frac{\lambda}{\pi w(z)^2}.
% \end{equation}

% Propagation through any ABCD system maps $q$ by the \textbf{M\"{o}bius
% (bilinear) transformation}
% \begin{equation}
%   q_{\text{out}} = \frac{A\,q_{\text{in}} + B}{C\,q_{\text{in}} + D}.
% \end{equation}
% This single complex equation encodes the full evolution of both the beam width
% and the wavefront curvature through an arbitrary paraxial system.

% \subsection*{Free-space propagation of a Gaussian beam}

% Using $M_{\text{drift}}$ with $A=D=1$, $B=d$, $C=0$:
% \begin{equation}
%   q_{\text{out}} = q_{\text{in}} + d,
% \end{equation}
% so the beam parameter simply \emph{advances} by $d$, in direct analogy with the
% shear of the ray cloud in phase space.

